\section*{Capítulo \ref{ch:g12:doppler-effect} --- O efeito Doppler}

\begin{solution}{exo:g12:doppler-effect:1}
(a) $f' = 700 \times 340/315 \approx \qty{756}{Hz}$.
(b) $f' = 700 \times 340/365 \approx \qty{652}{Hz}$.
(c) \qty{700}{Hz}: motorista e sirene se movem juntos, e a distância
entre eles nunca muda.
\end{solution}

\begin{solution}{exo:g12:doppler-effect:2}
A \emph{intensidade sonora} diminui porque a \omterm{def:g10:signals-and-waves:wave}{onda} se espalha --- ela
depende da distância. A \emph{altura} cai da aproximação para o
afastamento --- é o \omterm{def:g12:doppler-effect:doppler}{efeito Doppler}, que depende da velocidade com que a
distância muda. Dois efeitos independentes.
\end{solution}

\begin{solution}{exo:g12:doppler-effect:3}
Observador em movimento:
$f' = 700\,(1 + 6.0/340) \approx \qty{712}{Hz}$; pedalando para longe,
$700\,(1 - 6.0/340) \approx \qty{688}{Hz}$. A sirene está em repouso no
ar, de modo que o padrão de \omterm{def:g10:signals-and-waves:wave}{ondas} não é perturbado; só o observador corre
através dele.
\end{solution}

\begin{solution}{exo:g12:doppler-effect:4}
À frente: $\lambda' = 310/400 = \qty{0.775}{m}$; atrás:
$370/400 = \qty{0.925}{m}$; em repouso:
$340/400 = \qty{0.85}{m}$. O \omterm{def:g12:mechanical-waves:wavelength}{comprimento de onda} curto alcança quem o
trem se aproxima, e o longo, quem ele deixa para trás.
\end{solution}

\begin{solution}{exo:g12:doppler-effect:5}
$\Delta f/f \approx 5.0/340 \approx 1.5\%$: $\Delta f \approx
\qty{6.5}{Hz}$, ou seja, cerca de \qty{446}{Hz} --- um quarto de
semitom. Um músico perceberia, por pouco.
\end{solution}

\begin{solution}{exo:g12:doppler-effect:6}
A parede recebe $f\,v/(v - v_b)$ (fonte em movimento); o morcego,
observador em movimento que se aproxima do \omterm{rem:g10:signals-and-waves:honest}{eco}, escuta $(1 + v_b/v)$
vezes isso: $f' = f(v + v_b)/(v - v_b) = 50.0 \times 346/334 \approx
\qty{51.8}{kHz}$, um desvio de $+\qty{1.8}{kHz}$. É o deslocamento duplo
do \omterm{met:g10:signals-and-waves:echo}{radar} de trânsito com os papéis de emissor e refletor trocados: só a
velocidade de aproximação importa.
\end{solution}

\begin{solution}{exo:g12:doppler-effect:7}
$u = c\,\Delta f/(2f) = \num{3.00e8} \times 4000/(2 \times \num{24.15e9})
\approx \qty{24.8}{m/s} \approx \qty{89}{km/h}$. A \qty{50}{km/h}
($\qty{13.9}{m/s}$):
$\Delta f = 2 \times 13.9 \times \num{24.15e9}/\num{3.00e8} \approx
\qty{2.2}{kHz}$.
\end{solution}

\begin{solution}{exo:g12:doppler-effect:8}
$u = v\,\Delta f/(2f) = 1540 \times 2600/\num{1.0e7} \approx
\qty{0.40}{m/s}$. Um fluxo perpendicular tem \omterm{def:g12:doppler-effect:radial}{velocidade radial} nula, e
portanto desvio algum.
\end{solution}

\begin{solution}{exo:g12:doppler-effect:9}
(a) Fonte: $f' = 1000 \times 340/306 \approx \qty{1111}{Hz}$.
(b) Observador: $f' = 1000 \times 1.1 = \qty{1100}{Hz}$. A fórmula de
baixa velocidade prevê \qty{1100}{Hz}: exata para o observador e
\qty{11}{Hz} abaixo para a fonte --- é a correção de ordem $x^2$, $1\%$
de $f$ para $x = 0.1$.
\end{solution}

\begin{solution}{exo:g12:doppler-effect:10}
$f = 2 \times 465 \times 415/880 \approx \qty{439}{Hz}$;
$v_s = 340 \times 50/880 \approx \qty{19.3}{m/s} \approx \qty{70}{km/h}$.
O $f$ verdadeiro é a média harmônica: a aproximação eleva a \omterm{def:g10:signals-and-waves:frequency}{frequência}
mais do que o afastamento a reduz, e por isso $f$ fica abaixo da média
aritmética (\qty{440}{Hz}).
\end{solution}

\begin{solution}{exo:g12:doppler-effect:11}
\omterm{def:g12:mechanical-waves:wavelength}{Comprimento de onda} maior: \omterm{def:g12:doppler-effect:redshift}{desvio para o vermelho}, e a estrela se afasta.
$v_r = c\,\Delta\lambda/\lambda = \num{3.00e8} \times 0.2/656.3 \approx
\qty{9.1e4}{m/s} \approx \qty{91}{km/s}$.
\end{solution}

\begin{solution}{exo:g12:doppler-effect:12}
$\Delta\lambda = \lambda K/c = 500 \times 55/\num{3.00e8} \approx
\qty{9.2e-5}{nm}$;
$\Delta\lambda/\lambda = 55/\num{3.00e8} \approx \num{1.8e-7}$. O \omterm{def:g10:signals-and-waves:signal}{sinal}
inteiro são duas partes em dez milhões: um espectrógrafo que derive mais
do que $10^{-7}$ o soterra.
\end{solution}

\begin{solution}{exo:g12:doppler-effect:13}
$v_r = 0.024\,c \approx \qty{7.2e6}{m/s}$. Velocidades proporcionais à
distância, iguais em todas as direções: todas as distâncias se esticando
de modo uniforme --- um universo em expansão sem centro privilegiado. A
$2.4\%$, as correções são de ordem
$(v_r/c)^2 \approx \num{6e-4}$: ainda confiável.
\end{solution}

\begin{solution}{exo:g12:doppler-effect:14}
As duas bordas se deslocam em sentidos opostos, de modo que a linha média
do disco fica \emph{alargada}, e não deslocada: largura total
$2\lambda v/c = 2 \times 656.3 \times \num{2.0e3}/\num{3.00e8}
\approx \qty{8.8e-3}{nm}$ (cerca de $\pm\qty{4.4e-3}{nm}$).
\end{solution}

\begin{solution}{exo:g12:doppler-effect:15}
(a) À frente da viatura as cristas ficam espaçadas de $(v - v_s)T$; o
caminhão foge com $v_o$ e as cruza com velocidade relativa $v - v_o$:
$f' = (v - v_o)/\bigl((v - v_s)T\bigr) = f(v - v_o)/(v - v_s)$.
(b) $f' = 700 \times 310/300 \approx \qty{723}{Hz}$. (c) $v_o = v_s$ dá
$f' = f$: a separação é constante, e nenhuma mudança de distância
significa desvio Doppler nenhum, por definição.
\end{solution}

\begin{solution}{pb:g12:doppler-effect:1}
\textbf{1.} Cada crista se expande a partir do ponto em que foi emitida:
as cristas se amontoam à frente da buzina em movimento e se esticam atrás
dela, de modo que o celular que se aproxima as cruza mais depressa
($f_1$) do que o que se afasta ($f_2$).

\textbf{2.} $f_1 = f\,\dfrac{v}{v - u}$, \quad
$f_2 = f\,\dfrac{v}{v + u}$.

\textbf{3.} $f_1 - f_2 = \frac{2fv\,u}{v^2 - u^2}$ e
$f_1 + f_2 = \frac{2fv^2}{v^2 - u^2}$; dividindo,
$u = v\,\frac{f_1 - f_2}{f_1 + f_2} = 340 \times 60/882 \approx
\qty{23}{m/s} \approx \qty{83}{km/h}$.

\textbf{4.} $\dfrac{2f_1 f_2}{f_1 + f_2} = f$ (o fator comum se cancela):
$f = 2 \times 471 \times 411/882 \approx \qty{439}{Hz}$ --- a média
harmônica, abaixo da média \qty{441}{Hz} porque a aproximação eleva $f$
mais do que o afastamento o reduz.

\textbf{5.} À frente: $(v - u)/f = 316.9/439 \approx \qty{0.72}{m}$;
atrás: $363.1/439 \approx \qty{0.83}{m}$; em repouso:
$340/439 \approx \qty{0.77}{m}$.

\textbf{6.} Uma vez como \emph{observador} em movimento que recebe a
\omterm{def:g10:signals-and-waves:wave}{onda}, e outra como \emph{fonte} em movimento que reemite o que recebeu.

\textbf{7.} $f' = f\,(1 + u/c)/(1 - u/c) \approx f(1 + 2u/c)$, logo
$\Delta f = f' - f \approx \dfrac{2u}{c}\,f$.

\textbf{8.} $u = c\,\Delta f/(2f) = \num{3.00e8} \times
3550/(2 \times \num{24.125e9}) \approx \qty{22}{m/s}$.

\textbf{9.} $22.1 \times 3.6 \approx \qty{79}{km/h}$: logo abaixo do
limite de \qty{80}{km/h}.

\textbf{10.} $\Delta f = 2f/(3.6\,c) \approx \qty{45}{Hz}$ por
\unit{km/h}. Bater o \omterm{rem:g10:signals-and-waves:honest}{eco} contra a \omterm{def:g10:signals-and-waves:wave}{onda} emitida deixa apenas a
\emph{diferença}: uma \omterm{def:g10:signals-and-waves:frequency}{frequência} de áudio, contada trivialmente, embora
seja uma fração $10^{-7}$ da portadora.

\textbf{11.} Os \omterm{def:g12:mechanical-waves:wavelength}{comprimentos de onda} de \omterm{def:g10:light-spectra:emission}{linhas espectrais} conhecidas,
noite após noite; e depois $v_r = c\,\Delta\lambda/\lambda$.

\textbf{12.} $K = \qty{55}{m/s}$, $P = \qty{4.2}{d}$.

\textbf{13.} $\Delta\lambda = 550 \times 55/\num{3.00e8} \approx
\qty{1.0e-4}{nm}$, $\Delta\lambda/\lambda \approx \num{1.8e-7}$ --- uma
parte em cinco milhões, umas $\num{4e5}$ vezes menor do que o desvio de
$7\%$ do trem.

\textbf{14.} Estrela e planeta se atraem igualmente, de modo que os dois
orbitam o \omterm{def:g11:forces-and-motion:com}{centro de massa} comum: a estrela não pode ficar parada. Um
$v_r$ senoidal é um \omterm{def:g11:forces-and-motion:ucm}{movimento circular uniforme} visto de perfil --- uma
\omterm{def:g10:universal-gravitation:satellite}{órbita} circular.

\textbf{15.} Só a componente na linha de visada muda a distância
estrela--Terra e, portanto, o \omterm{def:g12:mechanical-waves:wavelength}{comprimento de onda}. De frente,
$v_r = 0$ o tempo todo: \omterm{def:g10:signals-and-waves:signal}{sinal} nenhum.

\textbf{16.} $P = \qty{3.63e5}{s}$:
$r^3 = \num{6.67e-11} \times \num{2.0e30} \times
(\num{3.63e5})^2/(4\pi^2) \approx \num{4.4e29}$, logo
$r \approx \qty{7.6e9}{m}$ --- cerca de $1/20$ da distância Terra--Sol.

\textbf{17.} $V = 2\pi r/P = 2\pi \times \num{7.6e9}/\num{3.63e5}
\approx \qty{1.3e5}{m/s}$.

\textbf{18.} $m = M K/V = \num{2.0e30} \times 55/\num{1.3e5} \approx
\qty{8.3e26}{kg}$.

\textbf{19.} Cerca de $0.4$ massa de Júpiter ($\approx 140$ Terras),
orbitando $8$ vezes mais perto do que Mercúrio: um gigante gasoso
assando contra a sua estrela --- um ``Júpiter quente''.

\textbf{20.} Uma inclinação esconde parte do movimento: o $K$ medido é só
a parcela radial, de modo que a massa verdadeira é \emph{maior} ---
\qty{8.3e26}{kg} é um mínimo. Um desvio \omterm{def:g10:signals-and-waves:period}{periódico} de uma parte em cinco
milhões acabou de entregar a \omterm{def:g10:universal-gravitation:satellite}{órbita}, a velocidade e a massa mínima de um
planeta que ninguém jamais viu.
\end{solution}
